\section{Results}

In this section we present results of this analysis. Results obtained with different measurement methods have different $p_T$ ranges some of which overlap. In order to merge different measurement methods results together we compared $K^*(892)$ invariant $p_T$ spectra results obtained with different methods which is shown in Figure \ref{fig:spectra_comp}. 

\begin{figure}[h!]
	\centering
	\foreach \centr in {0-20, 20-40, 40-60, 60-93, MB}
	{
		\begin{subfigure}{0.32\textwidth}
			\includegraphics[width=\textwidth]{InvM/Run7AuAu200/19079/Spectra_KStar892_\centr_comp.pdf}
		\end{subfigure}
	}
	\caption{Comparison of invariant $p_T$ spectra of $K^*(892)$ for different measurement methods in different centrality classes}
	\label{fig:spectra_comp}
\end{figure}

$K^*(892)$ invariant $p_T$ spectra are consistent within marker sizes therefore we decided to merge different methods results into one result. For overlapping $p_T$ bins for different pair mixing methods we chose the method that yields the smallest value of propagated statistical and systematic uncertainties. Propagated statistical and systematic uncertainties for every centrality class and for every pair mixing mehthod are presented in Appendix \ref{appendix:comp_uncertainties}. Table \ref{table:methods_usage} shows pair mixing methods with the smallest propagated statistical and systematic uncertaintites at every $p_T$ bin in every centrality class. 

To check the consistency of different measurement methods results within uncertainties the ratio of different measurements results to levy approximation of \Kstar invariant $p_T$ spectra (Section \ref{section:spectra}) were calculated and are shown on Figure \ref{fig:spectra_ratio}.

\begin{figure}[h!]
	\centering
	\foreach \centr in {0-20, 20-40, 40-60, 60-93, MB}
	{
		\begin{subfigure}{0.47\textwidth}
			\includegraphics[width=\textwidth]{InvM/Run7AuAu200/19079/Spectra_KStar892_\centr_ratio.pdf}
		\end{subfigure}
	}
	\caption{Comparison of ratio of invariant $p_T$ spectra of $K^*(892)$ to levy approximations for different measurement methods in different centrality classes. Error bars show propagated statistical and systematic uncertaintites.}
	\label{fig:spectra_ratio}
\end{figure}

\begin{table}[h!]
	\centering
	\begin{tabular}{c||c|c|c|c|c}
		$p_T$, GeV/c & 0-20\% & 20-40\% & 40-60\% & 60-93\% & MB \\
		\hline\hline
		0.9 - 1.1 & 2PID & 2PID & 2PID & 2PID & 2PID \\
		1.1 - 1.4 & 2PID & 2PID & 2PID & 2PID & 2PID \\
		1.4 - 1.7 & 2PID & 2PID & 2PID & 2PID & 2PID \\
		1.7 - 1.9 & 2PID & 2PID & 2PID & 2PID & 2PID \\
		1.9 - 2.1 & 2PID & 2PID & 2PID & 1PID & 2PID \\
		2.1 - 2.3 & 2PID & 2PID & 2PID & 1PID & 2PID \\
		2.3 - 2.6 & 2PID & 2PID & 2PID & 1PID & 2PID \\
		2.6 - 2.9 & 2PID & 1PID & 1PID & 1PID & 2PID \\
		2.9 - 3.4 & 1PID & 1PID & 1PID & 1PID & 1PID \\
		3.4 - 4 & 1PID & noPID & 1PID & noPID & 1PID \\
		4 - 4.5 & noPID & noPID & noPID & noPID & noPID \\
		4.5 - 5 & noPID & noPID & noPID & noPID & noPID \\
		5 - 6.5 & noPID & noPID & noPID & noPID & noPID \\
	\end{tabular}
	\caption{Pairs mixing methods that yield the lowest value of propagated statistical and systematic uncertainties for different centrality classes and for different $p_T$ bins}
	\label{table:methods_usage}
\end{table}

\subsection{\Kstar invariant $p_T$ spectra}
\label{section:spectra}

\Kstar invariant $p_T$ spectra was calculated by the following formula

\begin{equation}
   \frac{1}{2\pi} \frac{d^2N}{dy dp_T} = \frac{1}{2\pi p_T \epsilon} \frac{d^2 Y_{raw}}{dy dp_T}
\end{equation}

Where $Y_{raw}$ is a raw yield.

The result is shown on Figure \ref{fig:spectra}.

Invariant $p_T$ spectra was also approximated with exponential function (Eq. \ref{eq:exp_fit}) and with Levy function (Eq. \ref{eq:levy_fit}) \cite{levy1, levy2}. 
%Approximations parameters are shown in Tables \ref{table:exp_fit_parameters} - \ref{table:levy_fit_parameters}.

\begin{equation}
   \frac{1}{2\pi} \frac{d^2N}{dy dp_T} = \frac{1}{2 \pi} \frac{dN}{dy} \frac{1}{T^2} e^{- p_T/T}
   \label{eq:exp_fit}
\end{equation}

\begin{equation}
   \frac{1}{2\pi} \frac{d^2N}{dy dp_T} = \frac{1}{2\pi} \frac{dN}{dy} \frac{(n-1)(n-2)}{(\Lambda + m_{K^*(892)}(n-1))(\Lambda + m_{K^*(892)})} \left( \frac{\Lambda + \sqrt{p_T^2 + m_{K^*(892)}^2}}{\Lambda + m_{K^*(892)}} \right)^{-n}
   \label{eq:levy_fit}
\end{equation}

Where
\begin{itemize}
   \item $T$, $dN/dy$, $n$, $\Lambda$ - free parameters
   \item $m_{K^*(892)}$ - mass of \Kstar
\end{itemize}

\begin{figure}[h!]
	\centering
	\includegraphics[width=0.8\textwidth]{InvM/Run7AuAu200/19079/Spectra_KStar892.pdf}
	\caption{\Kstar invariant $p_T$ spectra. Dashed lines are exponential approximations, colored solid lines are Levy approximations.}
	\label{fig:spectra}
\end{figure}

\subsection{\Kstar nuclear modification factors}

\Kstar nuclear modification factors $R_{AA}$ were calculated by the following formula

\begin{equation}
	R_{AA} = \frac{1}{N_{coll}} \frac{d^2 N_{AA} / dy dp_T}{d^2 N_{pp} / dy dp_T}
\end{equation}

Where 
\begin{itemize}
	\item $N_{coll}$ is a Au+Au scaling factor or an average number of nucleon-nucleon collisions in Au+Au collision.
	\item $d^2 N_{AA} / dy dp_T$ - \Kstar invariant $p_T$ spectra in Au+Au
	\item $d^2 N_{pp} / dy dp_T$ - \Kstar invariant $p_T$ spectra in pp (AN911)
\end{itemize}

Values of $N_{coll}$ were taken from \cite{ncoll} and are presented on the Table \ref{table:ncoll}.

\begin{table}[h!]
	\centering
	\begin{tabular}{c|c}
		Centrality class, \% & $N_{coll}$ \\
		\hline
		0-20 & $784 \pm 77$ \\
		20-40 & $300.8 \pm 31.4$ \\
		40-60 & $94.2 \pm 13$ \\
		60-93 & $14.8 \pm 4$ \\
		MB & $257.8 \pm 25.4$ \\
	\end{tabular}
	\caption{$N_{coll}$ values for different centrality classes for Au+Au at \snn}
	\label{table:ncoll}
\end{table}

Nuclear modification factors $R_{AB}$ are shown on Figure \ref{fig:rab}.

\begin{figure}[h!]
	\centering
	\includegraphics[width=0.99\textwidth]{InvM/Run7AuAu200/19079/RAB_KStar892.pdf}
	\caption{Nuclear modification factors $R_{AB}$ of \Kstar, $\varphi(1020)$, $\pi^0$, and protons for different centrality classes and $p_T$ bins. Gray box on the right is a scaling uncertainty along y axis from $N_{coll}$.}
	\label{fig:rab}
\end{figure}

\newpage

Nuclear modification factors $R_{CP}$ were calculated using the following formula

\begin{equation}
	R_{CP} = \frac{N_{coll}^{peripheral}}{N_{coll}^{central}} \frac{d^2 N_{central} / dy dp_T}{d^2 N_{peripheral} / dy dp_T}
\end{equation}

Where 
\begin{itemize}
	\item $central$ denotes the values for central collision centrality class
	\item $peripheral$ denotes the values for peripheral collision centrality class
\end{itemize}

Nuclear modification factors $R_{CP}$ are presented on Figure \ref{fig:rcp}.

\begin{figure}[h!]
	\centering
	\includegraphics[width=0.99\textwidth]{InvM/Run7AuAu200/19079/RCP_KStar892.pdf}
	\caption{Nuclear modification factors $R_{CP}$ of \Kstar, charged pions, kaons, and protons. Grey and cyan boxes on the right of each picture show the scaling uncertainty along y axis from $N_{coll}$ of 0-20/60-93 and 0-10/60-93 respectively (left picture) and the uncertainty of $N_{coll}$ of 0-20/40-60 and 0-10/60-93 respectively (right picture).}
	\label{fig:rcp}
\end{figure}
